% Test-case specification TC-0020 (CTest test TC-0020-settings).
\documentclass[letterpaper]{article}

\usepackage[letterpaper,margin=1in]{geometry}
\emergencystretch=3em

% Compile from the repository root so this path resolves to
% the pinned WSP submodule.
\input{wsp/testing/test-case-library.tex}

\begin{document}

\def\TCID{TC-0020}
\def\TCTitle{%
  SaveSetting, GetSetting, GetAllSettings and DeleteSetting}
\def\TCPurpose{%
  Verify the registry-backed application settings functions.}
\def\TCPriority{\TCPriorityHigh}

\def\TCRequirementRef{%
  REQ-0004}
\def\TCDesignRef{%
  Not applicable.}
\def\TCFeatureRef{%
  \texttt{tests/\allowbreak{}vbrun\_\allowbreak{}settings\_\allowbreak{}tests.\allowbreak{}cpp},
  \texttt{src/\allowbreak{}helpers\_\allowbreak{}settings.\allowbreak{}cpp}}
\def\TCTestTechnique{%
  Requirements-based testing.}

\def\TCPreconditions{%
  The library and the test program are built
  (\texttt{scripts/\allowbreak{}Invoke-Build.\allowbreak{}ps1}) and the test
  is registered with CTest.}
\def\TCEnvironment{%
  Windows with the toolchain of
  \texttt{scripts/\allowbreak{}Invoke-Build.\allowbreak{}ps1}; the local
  gate builds and runs x64 Debug, CI builds and runs the other
  architectures. The library is built from this repository in
  \texttt{out/\allowbreak{}build}.}
\def\TCAssumptions{%
  Not applicable.}

\def\TCInputData{%
  A registry key of the test with the sections \texttt{Window} and
  \texttt{Colors}; an empty application name.}
\def\TCInitialState{%
  The test's registry key is absent.}

\def\TCProcedure{%
  \begin{enumerate}
    \item CTest runs \texttt{vbrun\_\allowbreak{}settings\_\allowbreak{}tests}.
    \item Save settings, read them back with \texttt{GetSetting} (including a
      missing setting and a default) and list them with
      \texttt{GetAllSettings}, a two-dimensional array of keys and values.
    \item Delete a setting, a section and the application and check that the
      others stay, that deleting a missing setting is error 5 and that an
      empty application name is error 5.
    \item CTest records the exit code and, on failure, the program's output
      (\texttt{--output-on-failure}).
  \end{enumerate}}
\def\TCExpectedResult{%
  All checks pass; the program prints \texttt{settings ok} and exits 0.}
\def\TCPostConditions{%
  The test's registry key is removed at the end.}

\TCTemplate

\end{document}
